% GLB
% Here we describe the basic API. This is now quite simple. user has
% to write a TaskFrame, and optionally a TaskBag, supply some
% parameters, and run the darn thing. We illustrate in running text 
% with a Fib example.

\subsection{The GLB Framework Design}
The GLB (Global Load-Balancing) framework is applicable for {\em
  tree-structured problems}. 

Such problems can be characterized thus. There is an initial
bag (multiset) of {\em tasks} (may be of size one). A task usually has a
small amount of associated state, but is permitted to access
(immutable) ``reference state''  that can be replicated across all
places. Consequently, tasks are {\em relocatable}: they can be
executed at any place. 

Tasks can be executed. Executing a task may only involve local,
self-contained work (no communication). During execution, a task is
permitted to create zero or more tasks, and produce a result of a given
(pre-specified) type \code{Z}. The user is required to specify a
reduction operator of type \code{Z}. 

The GLB framework is responsible for distributing the generated
collection of tasks across all nodes. Once no more tasks are left to
be executed, the single value of type \code{Z} (obtained by reducing
the results produced by all tasks) is returned.  

Since the execution of each task depends only on the state of the
task (and on immutable external state), the set of results produced on
execution will be the same from one run to another, regardless of
how the tasks are distributed across places. Since the user supplied
reduction operator is assumed to be associative and commutative, the
result of execution of the problem is determinate. Thus GLB is a
determinate application library that requires the user to provide a
few sequential pieces of code and handles concurrency, distribution
and communication automatically. 

The GLB framework is applicable to a wide variety of tasks. For a
simple example, consider the problem of computing the $n$'th Fibonacci
number in the naive recursive way. Here the state of the task can be
summarized by just one number (long), i.e.{} $n$. Execution of the
task yields either a result (for $n < 2$), or two tasks (for $n-1$ and
$n-2$). The results from all tasks need to be reduced through
arithmetic addition. 

All {\em state space search algorithms} from AI fall in this
category. Such algorithms are characterized by a space of states, an
initial state, and a generator function which given a state generates
zero or more states. The task is to enumerate all states (perhaps with
a cutoff), apply some function, and combine the results. An example of
such an application is the famous $N$-Queens problem.

Another example of this problem is the Unbalanced Tree Search (UTS)
benchmark. Indeed, it was our attempt to solve the UTS problem that
led to the creation of the GLB framework.

\subsection{The GLB API}
The GLB framework works by distributing bags of tasks from busy places
to idle places automatically. In special cases, the application
programmer may improve overall performance by specifying a custom
implementation of \code{TaskBag}. The implementation must support a
\code{split} operation which permits the current task bag to be split
into two (so the other piece can be distributed to another place) and
a \code{merge} operation which permits incoming work to be included in
the task bag at the current place. 

Thus in summary, the application programmer provides the following
pieces of code:
\begin{enumerate}
\item A \code{TaskFrame[Z]} class. This class is expected to maintain a
  reference to a \code{TaskBag} (the {\em local task bag}). The user
  must supply implementations for the following methods:
  \begin{enumerate}
  \item   \code{getTaskBag():TaskBag}: should return the local task bag.
  \item \code{getReducer():Reducible[Z]}: should return the reducer
    associated with this job.
  \item \code{getResult():Z}: This returns the accumulated result for all
    tasks that were completed at this place. 
  \item \code{initTask():void}: This is called once, when this task
    frame is set up at the current place. It should initialize the
    local task bag appropriately. 

    If the problem has a single root task, then the user should choose
    one place (e.g. place zero) as the place at which to add this
    task. If the problem has many initial tasks then more than one
    place may start out with non-empty task bags. 
  \item \code{runAtMostNTasks(n:Long):Boolean}: This method must run at most
    \code{n} tasks from the task bag at the current place, and return
    \code{true} if there is still work left to do in the local task
    bag. During the execution of this method, user code may manipulate
    the local task bag freely (removing and adding tasks to it). The
    framework guarantees that no other worker will concurrently access
    the task bag. For tasks which have completed and produced a result
    of type \code{Z}, the method must accumulate the result in local
    state, using the supplied reducer.
\end{enumerate}
\item Optionally, a custom class extending \code{TaskBag} and
  appropriately implementing \code{split} and \code{merge} methods. 
\end{enumerate}

The framework is invoked in a fairly simple way:

{\footnotesize
\begin{verbatim}
val r = new GlobalLoadBalancer[T](parameters);
val result = g.run(()=> new TaskFrameImp[T]());
\end{verbatim}}

\noindent First an instance of \code{GlobalLoadBalancer} is created
with a specification of parameters for controlling various aspects of
global load-balancing. Next the framework is invoked, supplying a
function that generates an instance of the user-specified
\code{TaskFrame} class, with the appropriate result type. This method
returns the result of running the user's problem, balanced across all
places. 

The code in Figure~\ref{table:fib} illustrates the use of this
framework with a simple Fibonacci computation. 

%% Add Fibonacci code in totality.

\begin{table}
{\footnotesize
\begin{verbatim}
package fib;

import glb.TaskBag;
import glb.TaskFrame;
import glb.GlobalLoadBalancer;
import glb.ArrayListTaskBag;
import glb.GLBParameters;

public class Fib(n:Long) {
    class FibFrame extends TaskFrame[Long] {
        val bag = new ArrayListTaskBag[Long]();
        val reducer = Reducible.SumReducer[Long]();
        var result:Long=0;
        public def getTaskBag()=bag;
        public def getReducer()=reducer;
        public def getResult()= result;
        public def initTask() {
            if (here.id==0) bag.bag().add(n);
        }
        public def runAtMostNTasks(n:Long) {
            val b = bag.bag();
            for (1..n) {
                if (b.size()<=0) break;
                val x = b.removeLast(); // constant time
                if (x < 2) result += x;
                else {
                    b.add(x-1); // constant time
                    b.add(x-2);
                }
            }
            return b.size()>0;
        }
    }
    public def run():Long {
        val g = new GlobalLoadBalancer[Long](GLBParameters.Default);
        return g.run(()=>new FibFrame());
    }
    static def fib(n:Long):Long=n<2? n: fib(n-1)+fib(n-2);
    public static def main(args:Rail[String]) {
        val N = args.size < 1 ? 10 : Long.parseLong(args(0));
        val result = new Fib(N).run();
        
        Console.OUT.println("fib(" + N + ")=" 
           + result + " (actual=" + fib(N)+")");
    }
}
\end{verbatim}}
\caption{Fibonacci, with GLB}\label{table:fib}
\end{table}

\subsection{Implementation}

GLB is essentially a library form of UTS. Hence the basic structure of
the code is the same as for UTS. 


\subsection{Validation}
We have implemented three benchmarks using this framework --
Fibonacci, UTS and BC. In each case, for various input sizes we
checked the result obtained for these benchmarks with results obtained
through alternative means (e.g. running the benchmarks directly), and
verified the results were the same.

\subsection{Performance} Performance is reported in the UTS-G
and BC-G sections, through the performance of those benchmarks
compared with UTS and BC, respectively.
